% !TEX root = calculus.tex

\chapter{FUNCTION}
\label{function}
{\parindent=0pt

\rdr Functions are widely used in elementary mathematics.

\athr Yes, of course. You are familiar with \emph{numerical functions}. Moreover, you have worked already with different numerical functions. Nevertheless, it will be worthwhile to dwell on the concept of the function. To begin with, what is your idea of a function?

\rdr As I understand it, a function is a certain correspondence between \emph{two variables}, for example, between $x$ and $y$. Or rather, it is a dependence of a variable $y$ on a variable $x$.

\athr What do you mean by a ``variable''?

\rdr It is a quantity which may assume different values.

\athr Can you explain what your understanding of the expression ``a quantity assumes a value'' is? What does it mean? And what are the reasons, in particular, that make a quantity to assume this or that value? Don't you feel that the very concept of a variable quantity (if you are going to use this concept) needs a definition?

\rdr O.K., what if I say: a function $y = f (x)$ symbolizes a dependence of $y$ on $x$, where $x$ and $y$ are numbers.

\athr I see that you decided to avoid referring to the concept of a \emph{variable quantity}. Assume that $x$ is a number and $y$ is also a number. But then explain, please, the meaning of the phrase ``a dependence between two numbers''.

\rdr But look, the words ``an independent variable'' and ``a dependent variable'' can be found in any textbook on mathematics.

\athr The concept of a variable is given in textbooks on mathematics after the definition of a function has been introduced.

\rdr It seems I have lost my way.

\athr Actually it is not all that difficult ``to construct'' an image of a numerical function. I mean \emph{image}, not \emph{mathematical definition} which we shall discuss later.

In fact, a numerical function may be pictured as a ``black box'' that \emph{generates a number at the output in response to a number at the input}. You put into this ``black box'' a number (shown by $x$ in \hyperref[fig-10]{Figure \ref{fig-10}}) and the "black box" outputs a new number ($y$ in \hyperref[fig-10]{Figure \ref{fig-10}}).
\begin{figure}[!h]
\centering
\begin{tikzpicture}[line cap=round,line join=round,>=triangle 45,x=0.5cm,y=.6cm]
\clip(-8,-3) rectangle (7,4);
\fill[line width=2.pt,color=DimGray,fill=DarkGray,fill opacity=0.25,rounded corners=8pt] (3.,2.) -- (-3.,2.) -- (-3.,-1.) -- (3.,-1.) -- cycle;
\filldraw[DimGray] (2.5,1.75) -- (-2.5,1.75) -- (-2.5,1.1) -- (2.5,1.1) -- cycle;

\draw [line width=.75pt,color=IndianRed] (6.,3.) circle (12pt);
\draw [line width=.75pt,color=IndianRed] (6.,3.) circle (8pt);
%\draw [line width=1.2pt,color=DarkGray] (3.,2.)-- (-3.,2.);
%\draw [line width=1.2pt,color=DarkGray] (-3.,-1.)-- (3.,-1.);
%\draw [line width=1.2pt,color=DarkGray] (-3.,-1.)-- (-3,2.);
%\draw [line width=1.2pt,color=DarkGray] (3.,-1.)-- (3.,2.);
\draw [line width=.75pt,color=IndianRed] (-6.,-2.) circle (12pt);
\draw [line width=.75pt,color=IndianRed] (-6.,-2.) circle (8pt);
\draw (-.1,0.7)node[color=black] {\footnotesize \textit{``Black Box''}};
\draw (-0.1,0.)node[color=black] {\footnotesize Working as a };
\draw (0.,-0.6)node[color=black] {\footnotesize Function};
\draw (6,3) node[color=black] {$x$};
\draw (6,2.4) node[color=black](one) {};
\draw (-6,-2) node[color=black] {$y$};
\draw (-5.4,-2) node[color=black] (two) {};
\draw (0,1.85) node[color=black](onea){};
\draw (0,-.9) node[color=black](twoa){};
\path (one) edge[-{latex},line width=2.5pt,in=90,out=270,color=DodgerBlue] (onea);
\path (twoa) edge[-{latex},line width=2.5pt,in=0,out=270,color=DodgerBlue] (two);
\end{tikzpicture}
%\includegraphics[width=0.55\textwidth]{figures/fig-10.pdf}
\caption{A numerical function as a black box.}
\label{fig-10}
\end{figure}
Consider, for example, the following function: 
\begin{equation*}%
y = 4x^{2} - 1
\end{equation*}
If the input is $x= 2$, the output is $y=15$; if the input is $x=3$, the output is $y= 35$; if the input is $x=10$, the output is $y = 399$, etc.

\rdr What does this ``black box'' look like? You have stressed that \hyperref[fig-10]{Figure \ref{fig-10}} is only symbolic.

\athr In this particular case it makes no difference. It does not influence the essence of the concept of a function. But a function can also be ``pictured'' like this:
\begin{equation*}%
4 \Box^{2} - 1
\end{equation*}
The square in this picture is a ``window'' where you input the numbers. Note that there may be more than one ``window''. For example,
\begin{equation*}%
\frac{4 \Box^{2} - 1}{|\Box| + 1}
\end{equation*}

\rdr Obviously, the function you have in mind is 
\begin{equation*}%
\frac{4 x^{2} - 1}{|x| + 1}
\end{equation*}

\athr Sure. In this case each specific value should be input into both ``windows'' simultaneously. ``Black box'' working as a function
 \hyperref[fig-10]{Figure \ref{fig-10}}.

By the way, it is always important to see such a ``window'' (or ``windows'') in a formula describing the function. Assume, for example, that one needs to pass from a function $y = f (x)$ to a function $y = f(x - 1)$ (on a graph of a function this transition corresponds to a displacement of the curve in the positive direction of the $x$-axis by 1). If you clearly understand the role of such a ``window'' (``windows''), you will simply replace in this ``window'' (these ``windows'') $x$ by $x - 1$. Such an operation is Illustrated by \hyperref[fig-11]{Figure \ref{fig-11}} which represents the following function
\begin{equation*}%
y = \frac{4 x^{2} - 1}{|x| + 1}
\end{equation*}
\begin{figure}[!h]
\centering
\begin{tikzpicture}[line cap=round,line join=round,>=triangle 45,rounded corners,x=1cm,y=1cm]
%\clip(-8,-3) rectangle (7,4);
  \draw (-3,2) [color=IndianRed]rectangle (-1,0) node[pos=.5,align=center,] {$\dfrac{4x^{2}-1}{|x|+ 1} $};
  \draw (-1.5,1.5) [color=black]rectangle (2,4) node[pos=.5,align=center] {\Large$\dfrac{4 \Box^{2}-1}{|\Box|+ 1} $};
  \draw (3.5,1.75) [color=IndianRed] rectangle (6.5,3.75) node[pos=.5,align=center] { $\dfrac{4(x-1)^{2}-1}{|x-1|+ 1} $};
  \draw (1,4.25) [color=DodgerBlue]rectangle (3,5) node[pos=.5,align=center]{ $|x-1| $};
  \draw (1,.5) [color=DodgerBlue]rectangle (3,1.25) node[pos=.5,align=center] { $|x-1| $};
\draw (1.1,4.65) node[color=black](onea) {};
\draw (1.1,.85) node[color=black](twoa) {};  
\draw (-0.2,3.2) node[color=black](one) {};
\draw (-.2,2.35) node[color=black] (two) {};
\draw (3.65,2.75) node[color=black](three) {};
\draw (1.9,2.75) node[color=black] (threea) {};
\path (one) edge[{latex}-,line width=2.5pt,in=180,out=90,color=DodgerBlue] (onea);
\path (twoa) edge[-{latex},line width=2.5pt,in=270,out=180,color=DodgerBlue] (two);
\path (three) edge[{latex}-,line width=3.5pt,in=0,out=180,color=DodgerBlue] (threea);
\end{tikzpicture}
%\includegraphics[width=0.9\textwidth]{figures/fig-11.pdf}
\caption{Change in a function from $f(x) \to f(x-1)$.}
\label{fig-11}
\end{figure}
Obviously, as a result of substitution of $x - 1$ for $x$ we arrive at a new function (new ``black box'') 
\begin{equation*}%
\frac{4 (\Box - 1)^{2} - 1}{|\Box -1 | + 1}, \quad y = \frac{4 (x - 1)^{2} - 1}{|x -1 | + 1}
\end{equation*}

\rdr I see. If, for example, we wanted to pass from $y = f (x)$ to $y = f \left( \frac{1}{x} \right)$, the function pictured in \fig{fig-11} would be transformed as follows:
\begin{equation*}%
y = \frac{\dfrac{4}{x^{2}} -1 }{\dfrac{1}{|x |} + 1}
\end{equation*}

\athr Correct. Now. try to find $y = f (x)$ if
\begin{equation*}%
2 f	\left( \dfrac{1}{x } \right) - f(x) = 3x
\end{equation*}

\rdr I am at a loss. 

\athr As a hint, I suggest replacing $x$ by $\dfrac{1}{x } $.

\rdr This yields
\begin{equation*}%
2 f (x)	- f \left( \dfrac{1}{x } \right) = \frac{3}{x}
\end{equation*}
Now it is clear. Together with the initial equation, the new equation forms a system of two equations for $f (x)$ and $f \left( \dfrac{1}{x } \right) $:
\begin{equation*}
\left.
\begin{split}
2f \left( \dfrac{1}{x } \right) - f(x) & = 3x \\
2 f(x) - f \left( \dfrac{1}{x } \right)  & = \frac{3}{x}
\end{split}
\right\}
\end{equation*}

By multiplying all the terms of the second equation by 2 and then adding them to the first equation, we obtain
\begin{equation*}%
 f (x) = x + \frac{2}{x}
\end{equation*}


\athr Perfectly true.

\rdr In connection with your comment about the numerical function as a ``black box'' generating a numerical output in response to a numerical input, I would like to ask whether other types of ``black boxes'' are possible in calculus. 

\begin{figure}[!h]
\centering
\begin{tikzpicture}[line cap=round,line join=round,>=triangle 45,rounded corners,x=1cm,y=1cm,scale=1]
%  \draw[step=.5cm,gray,very thin] (-5,-5) grid (5,5);

\clip(-7.1,-.1) rectangle (5.1,5.1);
\draw (-4,4) [color=IndianRed]rectangle (-2,5) node[pos=.5,align=center]{\small number};
\draw (-.5,4) [color=IndianRed]rectangle (1.5,5) node[pos=.5,align=center]{\small {numerical} \\[-5pt] \small{ function}};
\draw (3,4) [color=IndianRed]rectangle (5,5) node[pos=.5,align=center]{\small {numerical} \\[-5pt] \small{ function}};

  \draw (-5.5,2) [color=IndianRed]rectangle (-3.5,3) node[pos=.5,align=center] {\small {numerical} \\[-5pt] \small{ function}};    
  \draw (-2,2) [color=IndianRed]rectangle (0,3) node[pos=.5,align=center] {\small operator};  
  \draw (1.5,2) [color=IndianRed]rectangle (3.5,3) node[pos=.5,align=center] {\small functional};

        \draw (-7,0) [color=IndianRed]rectangle (-5,1) node[pos=.5,align=center] {\small  number};    
    \draw (-3.5,0) [color=IndianRed]rectangle (-1.5,1) node[pos=.5,align=center] {\small {numerical} \\[-5pt] \small{ function}};
 \draw (0,0) [color=IndianRed]rectangle (2,1) node[pos=.5,align=center] {\small number};

\draw (4,4.1) node[color=black](1a) {};
\draw (2.5,2.8) node[color=black](1b) {};
\draw (2.5,2.1) node[color=black](1c) {};
\draw (1,.8) node[color=black](1d) {};

\draw (0.5,4.1) node[color=black](2a) {};
\draw (-1,2.8) node[color=black](2b) {};
\draw (-1,2.1) node[color=black](2c) {};
\draw (-2.5,.8) node[color=black](2d) {};

\draw (-3,4.1) node[color=black](3a) {};
\draw (-4.5,2.8) node[color=black](3b) {};
\draw (-4.5,2.1) node[color=black](3c) {};
\draw (-6,.8) node[color=black](3d) {};

\path (1a) edge[-{latex},line width=1.5pt,in=90,out=270,color=DodgerBlue] (1b);
\path (1c) edge[-{latex},line width=1.5pt,in=90,out=270,color=DodgerBlue] (1d);

\path (2a) edge[-{latex},line width=1.5pt,in=90,out=270,color=DodgerBlue] (2b);
\path (2c) edge[-{latex},line width=1.5pt,in=90,out=270,color=DodgerBlue] (2d);

\path (3a) edge[-{latex},line width=1.5pt,in=90,out=270,color=DodgerBlue] (3b);
\path (3c) edge[-{latex},line width=1.5pt,in=90,out=270,color=DodgerBlue] (3d);

\end{tikzpicture}
%\includegraphics[width=\textwidth]{figures/fig-12.pdf}
\caption{Understanding place and role of the numerical function as a mathematical tool.}
\label{fig-12}
\end{figure}

\athr Yes, they are. In addition to the numerical function, we shall discuss the concepts of an operator and a functional. 

\rdr I must confess I have never heard of such concepts. 

\athr I can imagine. I think, however, that \fig{fig-12} will be helpful. Besides, it will elucidate the place and role of the numerical function as a mathematical tool. \fig{fig-12} shows that:
\begin{enumerate}[label=$\textcolor{IndianRed}{\blacktriangleright}$]

\item \emph{a numerical function is a \emph{``black box''} that generates a number at the output in response to a number at the input;}

\item \emph{an operator	is	a \emph{``black box''} that	generates	a	numerical function at the output in response to a numerical function at the input}; it is said that all operator applied to a function generates a new function;

\item \emph{a functional is a \emph{``black box''} that generates a number at the output in response to a numerical junction at the input}, i.e. a concrete number is obtained ``in response'' to a concrete function.
\end{enumerate}

\rdr Could you give examples of operators and functionals?

\athr Wait a minute. In the next dialogues we shall analyze both the concepts of an operator and a functional. So far, we shall confine ourselves to a general analysis of both concepts. Now we get back to our main object, the numerical function.

The question is: How to construct a ``black box'' that generates a numerical function.

\rdr Well, obviously, we should find a relationship, or a law, according to which the number at the ``output'' of the ``black box'' could be forecast for each specific number introduced at the ``input''.

\athr You have put it quite clearly. Note that such a law could be naturally referred to as the \emph{law of numerical correspondence}. However, the law of numerical correspondence would not be a sufficient definition of a numerical function.

\rdr What else do we need?

\athr Do you think that any number could be fed into a specific ``black box'' (function)?

\rdr I see. I have to define a set of numbers acceptable as inputs of the given function.

\athr That's right. This set is said to be the \emph{domain of a function}.

Thus, the definition of a numerical function is based on two ``cornerstones'';

the domain of a function (a certain set of numbers), and the law of numerical correspondence. 

According to this law, \emph{every number from the domain of
a function is placed in correspondence with a certain number, which is called the value of the function; the values form the range of the function.}

\rdr Thus, we actually have to deal with two numerical sets. On the one hand, we have a set called the domain of a function and, on  the other, we have a set called the range of a function.

\athr At this juncture we have come closest to a mathematical definition of a function which will enable us to avoid the somewhat mysterious word ``black box''.

Look at \fig{fig-13}. It shows the function $y = \sqrt{1 - x^{2}}$.  \fig{fig-13} pictures two numerical sets, namely, $D$ (represented by the interval $[-1, 1]$) and	 $E$ (the interval $[0, 1]$). For your convenience these sets are shown on two different real lines.

\begin{figure}[!h]
\centering
\begin{tikzpicture}[line cap=round,line join=round,>=triangle 45,rounded corners,scale=1.5]
\clip(-2.2,-1.1) rectangle (4.5,5);
% \draw[step=.5cm,gray,very thin] (-3,-3) grid (6,6);
{\footnotesize
\draw (0,0) node[color=black,below](0a){0};
\draw (0,-.75) node[color=black,below]{D};
\draw (-2,0) node[color=black,below](-2a) {-1};
\draw (2,0) node[color=black,below](2a) {1};
\draw (-1.5,0) node[color=black,below](-15a) {};
\draw (1.5,0) node[color=black,below](15a) {};
\draw (-1,0) node[color=black,below](-1a) {};
\draw (1,0) node[color=black,below](1a) {};
\draw (-.5,0) node[color=black,below](-5a) {};
\draw (.5,0) node[color=black,below](5a) {};

\draw (3.5,1.5) node[color=black,right](0b){0};
\draw (4.2,2.7) node[color=black,right]{E};
\draw (3.5,4) node[color=black,right](2b) {1};
\draw (3.5,3.90) node[color=black,right](3b) {};% {0.96};
\draw (3.5,3.65) node[color=black,right](4b){};% {0.86};
\draw (3.5,3.15) node[color=black,right](5b){};% {0.66};


\draw [->,line width=1.2pt,color=DarkGray] (-2.5,0.) -- (3.5,0.);
\draw [->,line width=1.2pt,color=DarkGray] (3.5,1.) -- (3.5,5.);

\path (0a) edge[-{latex},line width=1.5pt,in=180,out=90,color=DodgerBlue] (2b);

\path (2a) edge[-{latex},line width=1.5pt,in=180,out=90,color=DodgerBlue] (0b);
\path (-2a) edge[-{latex},line width=1.5pt,in=180,out=90,color=DodgerBlue,looseness=.8] (0b);


\path (15a) edge[-{latex},line width=1.5pt,in=180,out=90,color=DodgerBlue] (5b);
\path (-15a) edge[-{latex},line width=1.5pt,in=180,out=90,color=DodgerBlue,looseness=.8] (5b);

\path (1a) edge[-{latex},line width=1.5pt,in=180,out=90,color=DodgerBlue] (4b);
\path (-1a) edge[-{latex},line width=1.5pt,in=180,out=90,color=DodgerBlue,looseness=.8] (4b);

\path (5a) edge[-{latex},line width=1.5pt,in=180,out=90,color=DodgerBlue] (3b);
\path (-5a) edge[-{latex},line width=1.5pt,in=180,out=90,color=DodgerBlue,looseness=.8] (3b);

\draw [decoration={brace,amplitude=20pt},decorate,rotate=-90,color=DarkGray,line width=.5pt] ($(2a)+(0.25em,.5ex)$) -- ($(-2a)+(0.25em,-.5ex)$);

\draw [decoration={brace,amplitude=20pt},decorate,rotate=-90,color=DarkGray,line width=.5pt] ($(2b)+(0.25em,.5ex)$) -- ($(0b)+(0.25em,.5ex)$);
}
\end{tikzpicture}

%\includegraphics[width=0.55\textwidth]{figures/fig-13.pdf}
\caption{Range and domain of a function.}
\label{fig-13}
\end{figure}

The set $D$ is the domain of the function, and $E$ is its range. Each number in $D$ corresponds to one number in $E$ (every input value is placed in correspondence with one output value). This correspondence is shown in \fig{fig-13} by arrows pointing from $D$ to $E$.


\rdr But \fig{fig-13} shows that two \emph{different} numbers in $D$ correspond to one number in $E$.

\athr It does not contradict the statement ``each number in $D$ corresponds to one number in $E$.'' I never said that \emph{different} numbers in $D$ must correspond to different numbers in $E$. Your remark (which actually stems from specific characteristics of the chosen function) is of no principal significance. Several numbers in $D$ may correspond to one number in $E$. An inverse situation, however, is forbidden. It is not allowed for one number in $D$ to correspond to more than one number in $E$. I emphasize that each number in $D$ must correspond to \emph{only one} (not more!) number in $E$.

Now we can formulate a mathematical definition of the numerical function.

\begin{mytheo}{Definition}
Take two numerical sets $D$ and $E$ in which each element $x$ of $D$ (this is denoted by $x \in D$) is placed in one-to-one correspondence with one element $y$ of $E$. Then we say that a function $y = f (x)$ is set in the domain $D$, the range of the function being $E$. It is said that the argument $x$ of the function $y$ passes through $D$ and the values of $y$ belong to $E$.
\end{mytheo}

Sometimes it is mentioned (but more often omitted altogether) that both $D$ and $E$ are subsets of the set of real numbers $R$ (by definition, $R$ is the real line).

On the other hand, the definition of the function can be reformulated using the term ``mapping''. Let us return again to \fig{fig-13}. Assume that the number of arrows from the points of $D$ to the points of $E$ is infinite (just imagine that such arrows have been drawn from each point of $D$). Would you agree that such a picture brings about an idea that $D$ is \emph{mapped} onto $E$?

\rdr Really, it looks like mapping.

\athr Indeed, this \emph{mapping} can be used to define the \emph{function}.
\begin{mytheo}{Definition}
A numerical function is a mapping of a numerical set $D$ (which is the domain of the function) onto another numerical set $E$ (the range of this function).
\end{mytheo}
Thus, the numerical function is a \emph{mapping of one numerical set onto another numerical set.} The term ``mapping'' should be understood as a kind of numerical correspondence discussed above. In the notation $y = f (x)$, symbol $f$ means the function itself (i.e. the mapping), with $x \in D$ and $y \in E$.

\rdr If the \emph{numerical function} is a mapping of one numerical set onto another numerical set, then the operator can be considered as a mapping of a set of numerical function onto another set of functions, and the \emph{functional} as a mapping of a set of functions onto a numerical set.

\athr You are quite right.

\rdr I have noticed that you persistently use the term ``numerical function'' (and I follow suit), but usually one simply says ``function''. Just how necessary is the word ``numerical''?

\athr You have touched upon a very important aspect. The point is that in modern mathematics the concept of a function is substantially broader than the concept of a numerical function. As a matter of fact, the concept of a function includes, as particular cases, a numerical function as well as an operator and a functional, because the essence in all the three is a mapping of one set onto another independently of the nature of the sets. You have noticed that both operators and functionals are mappings of certain sets onto certain sets. In a particular case of mapping of a numerical set onto a numerical set we come to a \emph{numerical} function. In a more general case, however, sets to be mapped can be \emph{arbitrary}. Consider a few examples.

\textcolor{IndianRed}{\textbf{Example 1}} Let $D$ be a set of working days in an academic year, and $E$ a set of students in a class. Using these sets, we can define a function realizing a schedule for the students on duty in the classroom. In compiling the schedule, each element of $D$ (every working day in the year) is placed in one-to-one correspondence with a certain element of $E$ (a certain student). This function is a \emph{mapping of the set of working days onto the set of students}. We may add that the domain of the function consists of the working days and the range is defined by the set of the students.

\rdr It sounds a bit strange. Moreover, these sets have finite numbers of elements.

\athr This last feature is not principal.

\rdr The phrase ``the values assumed on the set of students'' sounds somewhat awkward.

\athr Because you are used to interpret ``value'' as ``numerical value''.
Let us consider some other examples.

\textcolor{IndianRed}{\textbf{Example 2}} Let $D$ be a set of all triangles, and $E$ a set of positive real numbers. Using these sets, we can define two functions, namely, the area of a triangle and the perimeter of a triangle. Both functions are mappings (certainly, of different nature) \emph{of the set of the triangles onto the set of the positive real numbers}. It is said that the set of all the triangles is the domain of these functions and the set of the positive real numbers is the range of these functions.

\textcolor{IndianRed}{\textbf{Example 3}} Let $D$ be a set of all triangles, and $E$ a set of all circles. The mapping of $D$ onto $E$ can be either a circle inscribed in a triangle, or a circle circumscribed around a triangle. Both have the set of all the triangles as the domain of the function and the set of all the circles as the range of the function.

By the way, do you think that it is possible to ``construct'' an \emph{inverse} function in a similar way, namely, to define a function with all the circles as its domain and all the triangles as its range?

\rdr I see no objections.

\athr No, it is impossible. Because any number of different triangles can be inscribed in or circumscribed around a circle. In other words, each element of $E$ (each circle) corresponds to an Infinite number of different elements of $D$ (i.e, an infinite number of triangles). It means that there is no function since no mapping can be realized.

However, the situation can be improved if we restrict the set of triangles.

\rdr I guess I know how to do it. We must choose the set of all the \emph{equilateral} triangles as the set $D$. Then it becomes possible to realize both a mapping of $D$ onto $E$ (onto the set of all the circles) and an inverse mapping, i.e. the mapping of $E$ onto $D$, since only one equilateral triangle could be inscribed in or circumscribed around a given circle.

\athr Very good. I see that, you have grasped the essence of the concept of functional relationship. I should emphasize that from the broadest point of view this concept is based on the idea of mapping one set of objects onto another set of objects. \emph{It means that a function can be realized as a numerical function, an operator, or a functional}. As we have established above, a \emph{function may be represented by an area or perimeter of a geometrical figure, such, as a circle inscribed in a triangle or circumscribed around it, or it may take the form of a schedule of students on duty in a classroom, etc}. It is obvious that a list of different functions may be unlimited.

\rdr I must admit that such a broad interpretation of the concept of a function is very new to me.

\athr As a matter of fact, in a very diverse set of possible functions (mappings), we shall use only \emph{numerical functions, operators}, and \emph{functionals}. Consequently, we shall refer to numerical functions as simply \emph{functions}, while \emph{operators} and \emph{functionals} will be pointed out specifically.

And now we shall examine the already familiar concept of a numerical sequence as an example of mapping.

\rdr A numerical sequence is, apparently, a mapping of a set of natural numbers onto a different numerical set. The elements of the second set are the terms of the sequence. Hence, a numerical sequence is a particular case of a numerical function. The domain of a function is represented by a set of natural numbers.

\athr This is correct. But you should bear in mind that later on we shall deal with numerical functions whose domain is represented by the \emph{real line}, or by its \emph{interval} (or \emph{intervals}), and whenever we mention a function, we shall imply a numerical function.

In this connection it is worthwhile to remind you of the classification of intervals. In the previous dialogue we have already used this classification, if only partially.

First of all we should distinguish between the intervals of finite length:
\begin{enumerate}[label=$\textcolor{IndianRed}{\blacktriangleright}$]

\item a \emph{closed interval} that begins at $a$ and ends at $b$ is denoted by $[a, b]$; the numbers $x$ composing this interval meet the inequalities $a \leqslant x \leqslant b$;

\item an \emph{open interval} that begins at $a$ and ends at $b$ is denoted by $]\, a, b[$; the numbers $x$ composing this interval meet the inequalities	$a < x < b$;

\item a \emph{half-open interval} is denoted either by $] \, a, b]$ or $[a, b[$, the former implies that $ a < x \leqslant b$, and the latter that $a \leqslant x < b$.
\end{enumerate}

The intervals may also be \emph{infinite}:
\begin{center}
\begin{tabular}{llc}
$]\, - \infty, \, \infty \, [$ &	$( -	\infty < x < \infty)$ &	the real line \\
$] \, a, \, \infty  \, ]$ & $(a < x < \infty)$ & \\
$[\, a, \, \infty \, [$ &  $(a  \leqslant x < \infty)$ & \\
$] \, - \infty, b \, [$ & $(- \infty < x < b)$ & \\
$ ]- \infty, \, b]$  & $(-  \infty < x \leqslant b)$ & \\
\end{tabular}
\end{center}

Let us consider several specific examples of numerical functions. Judging by the appearance of the formulas given below, point out the intervals constituting the domains of the following functions:
\begin{align}%
y & = \sqrt{ 1 - x^{2}}	\label{fn-ex-01}\\
%eq(1) 
y & = \sqrt{ x - 1}		\label{fn-ex-02}\\
%eq(2) 
y & = \sqrt{2 - x} 	\label{fn-ex-03}\\
%eq(3)
y & = \dfrac{1}{\sqrt{x-1}}	\label{fn-ex-04}\\	
%eq(4)
y & = \dfrac{1}{\sqrt{2-x}}		\label{fn-ex-05}\\
%eq(5) 
y & = 	\sqrt{x-1} + \sqrt{2-x}\label{fn-ex-06}\\
%eq(6) 
y & = \dfrac{1}{\sqrt{x-1}}	+ \dfrac{1}{\sqrt{2 -x }}		\label{fn-ex-07}\\
%eq(7)
y & =	\sqrt{2 -x } +  \dfrac{1}{\sqrt{x-1}}	 	\label{fn-ex-08}\\
%eq(8)
y & = \sqrt{x-1} +	\dfrac{1}{\sqrt{2 -x }} \label{fn-ex-09}
%eq(9)
\end{align}
\rdr It is not difficult. The domain of function \eqref{fn-ex-01} is the interval $(-1, \, 1]$; that of \eqref{fn-ex-02} is $[1, \, \infty[$; that of \eqref{fn-ex-03} is $]-\infty, \, 2]$; that of \eqref{fn-ex-04} is $]\, 1, \infty[$; that of \eqref{fn-ex-05} is $]\infty, 2[$; that of \eqref{fn-ex-06} is $[1, 2]$, etc.

\athr Yes, quite right, but may I interrupt you to emphasize that if a function is a sum (a difference, or a product) of two functions, its domain is represented by the \emph{intersection} of the sets which are the domains of the constituent functions. It is well illustrated by function \eqref{fn-ex-06}. As a matter of fact, the same rule must be applied to functions \eqref{fn-ex-07}-\eqref{fn-ex-09}. Please, continue.

\rdr The domains of the remaining functions are \eqref{fn-ex-07} $]1, 2[$; \eqref{fn-ex-08} $]1,2]$; \eqref{fn-ex-09} $[1, 2[$.

\athr And what can you say about the domain of the function $y= \sqrt{x -2 } + \sqrt{1 -x }$?

\rdr The domain of  $y= \sqrt{x -2 }$ is $[2, \infty[$, while that of  $y= \sqrt{1 -x }$is $]- \infty, 1]$. These intervals do not intersect.

\athr It means	that the formula  $y= \sqrt{x -2 } + \sqrt{1 -x }$ does not define any function.
}